\chapter{Bit stuffing and CRC-15}
\label{ch:stuffing}

{\large\itshape The two mechanisms that let a frame survive a noisy bus.\par}

\begin{chapterbox}{In this chapter}
\begin{itemize}
  \item Why a long run of equal bits is a problem, and the rule of five that solves it.
  \item Which fields are stuffed, and which are not.
  \item CRC-15: the polynomial, the bit-serial engine, and an example worked by hand.
  \item Why the same engine both computes and checks.
\end{itemize}
\end{chapterbox}

\section{The problem with equal bits}

CAN nodes have no shared clock. Every node has its own oscillator, and they never run at exactly
the same speed. A receiver keeps in step by watching the \emph{edges} in the bit stream: every time
the signal changes level, it knows where a bit boundary was, and can adjust its own counter
(\kapref{ch:tajming}).

If no edges come, the receiver drifts. After enough equal bits in a row it samples the wrong bit,
and the frame is lost.

\textbf{The solution is to make the edges mandatory.} After five equal bits in a row, the
transmitter inserts a \textbf{stuff bit} of the opposite value, which the receiver recognises and
discards.

\section{The rule of five}

The rule is short: \textbf{after five equal bits in a row, a bit of the opposite value is sent},
whatever comes next.

An example, with the inserted bits marked:

\begin{console}
in:  1 1 1 1 1 0 0 0 0 0 1
out: 1 1 1 1 1 0 0 0 0 0 1 0 1
               ^         ^
               stuff     stuff
\end{console}

There are two details here that are easy to get wrong, and both are classic bugs:

\textbf{The stuff bit counts itself.} After a stuff bit has been inserted, the count starts again
with the stuff bit as the first bit of the new run. It is an entirely ordinary bit in the stream,
for everyone except whoever has to remove it again.

\textbf{The rule applies to the stream, not the field.} Five equal bits spanning the boundary
between the identifier and the DLC trigger a stuff bit just as five within one field do. An
implementation that resets the counter at every field change almost always works, and breaks
exactly when two fields happen to end and begin with the same bit value.

\subsection{Which fields are stuffed}

Stuffing applies from SOF up to and including the CRC field. After that it \textbf{stops}:

\begin{narrowtable}{@{}ll@{}}
\thead{Field} & \thead{Stuffed?} \\ \midrule
SOF, arbitration, control, data, CRC & Yes \\
CRC delimiter, ACK, EOF, intermission & No \\
\end{narrowtable}

That the end is not stuffed is what makes EOF usable as an end marker. Seven recessive bits in a
row cannot occur in the stuffed part --- there, five equal is the most that gets through --- so a
receiver that sees seven recessive knows for certain that the frame is over.

\subsection{Stuff errors}

The receiver counts the same way as the transmitter. After five equal bits, the next bit
\emph{must} be the opposite one. If it is not, something is wrong, and the receiver reports a
\textbf{stuff error}.

It is one of CAN's cheapest means of error detection: it costs no extra information at all, only a
counter at each end.

\begin{aside}
\note[The cost] Stuffing makes the frame longer, and by how much depends on the data. A frame whose
data bits alternate all the time gets no stuff bits at all; a frame full of zeros gets one every
fifth bit. For the longest data frame the overhead is at most twenty-odd bits. That is also why a
CAN frame's length cannot be predicted exactly from the DLC --- which in turn is why anyone
calculating bus load calculates with the worst case.
\end{aside}

\section{CRC-15}

The checksum is fifteen bits, computed over everything from SOF up to and including the data field
--- that is, the \emph{real} bits, without stuff bits. Stuff bits are the transport's business and
are not part of what is protected.

\subsection{The polynomial}

CAN uses

\[ x^{15} + x^{14} + x^{10} + x^{8} + x^{7} + x^{4} + x^{3} + 1 \]

which in binary form, the coefficients for $x^{14}$ down to $x^{0}$, becomes

\begin{console}
100010110011001
\end{console}

The polynomial is chosen to detect all single, double and triple errors, all error bursts up to
15 bits, and the vast majority of longer ones. It is not a random constant.

\subsection{The engine}

The computation is a fifteen-bit shift register with feedback, and it takes one bit at a time:

\begin{cppcode}
// One bit into the CRC-15 engine. crc is 15 bits wide.
const bool feedback{static_cast<bool>((crc >> 14U) & 1U) ^ bit};
crc = static_cast<std::uint16_t>((crc << 1U) & 0x7FFFU);
if (feedback)
{
    crc ^= 0x4599U;   // 100010110011001
}
\end{cppcode}

Three lines: XOR the incoming bit with the register's top bit, shift the register one step, and
XOR in the polynomial if the result of the first operation was one.

\subsection{An example by hand}

The sequence \code{10110010}, fed in bit by bit into a cleared register:

\begin{narrowtable}{@{}llll@{}}
\thead{Bit} & \thead{Value} & \thead{Register after} & \thead{Hex} \\ \midrule
1 & 1 & \code{100010110011001} & \code{0x4599} \\
2 & 0 & \code{100111010101011} & \code{0x4EAB} \\
3 & 1 & \code{001110101010110} & \code{0x1D56} \\
4 & 1 & \code{111111100110101} & \code{0x7F35} \\
5 & 0 & \code{011101111110011} & \code{0x3BF3} \\
6 & 0 & \code{111011111100110} & \code{0x77E6} \\
7 & 1 & \code{110111111001100} & \code{0x6FCC} \\
8 & 0 & \code{001101000000001} & \code{0x1A01} \\
\end{narrowtable}

The checksum is therefore \code{0x1A01}.

\subsection{The same engine checks}

This is where the elegance lies. The receiver has no separate checking function: it feeds the
\emph{same} bits plus \emph{the received checksum} through the same engine.

If everything is correct, the register reads \textbf{zero} once the last CRC bit has passed. If we
feed in \code{10110010} followed by \code{001101000000001} from the example above, the final value
is \code{0x0000}.

That is a property of the construction, not a special case: appending the remainder to the dividend
makes it divide evenly. Transmitter and receiver can therefore use exactly the same module, with no
mode switch --- which is exactly what a hardware design wants.

\begin{aside}
\note[A detail that bites] The engine still has to be resettable. An aborted frame leaves garbage in
the register, and the next frame that starts with a dirty register gets the wrong checksum from its
first bit. The reset happens at the start of every frame.
\end{aside}

\subsection{What the CRC does not protect}

The CRC covers SOF up to and including the data field. It therefore does \emph{not} cover the CRC
delimiter, the ACK field or EOF. Errors in those are detected by the format rules instead: a
delimiter that is not recessive is a form error (\kapref{ch:fel}).

\section*{Exercises}

\exercise{Stuff a stream}{Calculation}
\label{ex:3:1}
Stuff the following bit streams according to the rule of five, and mark the inserted bits.
\begin{enumerate}[a)]
  \item \code{0000011111}
  \item \code{1111111111}
  \item \code{0101010101}
\end{enumerate}
How many bits did each one become after stuffing?

\exercise{Destuff}{Calculation}
\label{ex:3:2}
A receiver receives \code{111110000010}.
\begin{enumerate}[a)]
  \item Which bits are stuff bits?
  \item What was the original bit stream?
  \item Suppose instead that the receiver got \code{111111...}. What does it report, and after
    which bit?
\end{enumerate}

\exercise{The CRC engine}{Calculation}
\label{ex:3:3}
Feed \code{1101} into a cleared CRC-15 register, one step at a time, and write down the register's
value after each bit. Use the table in section 3.3 as a template.

\exercise{Why zero?}{Reasoning}
\label{ex:3:4}
\begin{enumerate}[a)]
  \item Why does the receiver's register read zero once both the data and the checksum have been fed
    in?
  \item What practical advantage does that give a hardware design?
  \item A node by mistake feeds the \emph{stuff bits} into the CRC engine as well. What happens, and
    who detects it?
\end{enumerate}
